%% kapa: exam=1 round=34 year=2023 period=1 form=A subject=civil_law subject_ko=민법 seq=26 number=26
%% source: Q-Net 감정평가사 기출문제 (cst00309) artlSeq=5211789 file=raw/1008/5211789/2023년도 제34회 감정평가사 1차시험 1교시.pdf
%% conversion: kordoc 4.24.51 -> tools/gen1.py (auto); review=proofread-sonnet; flags=none
\begin{question}{26}
자주점유에 관한 설명으로 옳지 않은 것은? (다툼이 있으면 판례에 따름)
\begin{choices}
  \item 점유매개자의 점유는 타주점유에 해당한다.
  \item 부동산의 매매 당시에는 그 무효를 알지 못하였으나 이후 매매가 무효임이 밝혀지더라도 특별한 사정이 없는 한, 매수인의 점유는 여전히 자주점유이다.
  \item 양자간 등기명의신탁에 있어서 부동산 명의수탁자의 상속인에 의한 점유는 특별한 사정이 없는 한, 자주점유에 해당하지 않는다.
  \item 공유토지 전부를 공유자 1인이 점유하고 있는 경우, 특별한 사정이 없는 한 다른 공유자의 지분비율 범위에 대해서는 타주점유에 해당한다.
  \item 자주점유의 판단기준인 소유의 의사 유무는 점유취득의 원인이 된 권원의 성질이 아니라 점유자의 내심의 의사에 따라 결정된다.
\end{choices}
\end{question}
\answer{5}
